\documentclass[
	ngerman,
	twoside,
	pdfa=false,
	ruledheaders=section,%Ebene bis zu der die Überschriften mit Linien abgetrennt werden, vgl. DEMO-TUDaPub
	class=report,% Basisdokumentenklasse. Wählt die Korrespondierende KOMA-Script Klasse
	thesis={type=Hausarbeit},%bachelor},% Dokumententyp Thesis, für Dissertationen siehe die Demo-Datei DEMO-TUDaPhd
	accentcolor=2c,%3d,% Auswahl der Akzentfarbe
	custommargins=false,% Ränder werden mithilfe von typearea automatisch berechnet
	marginpar=false,% Kopfzeile und Fußzeile erstrecken sich nicht über die Randnotizspalte
	%BCOR=5mm,%Bindekorrektur, falls notwendig
	parskip=half-,%Absatzkennzeichnung durch Abstand vgl. KOMA-Sript
	fontsize=12pt,%Basisschriftgröße laut Corporate Design ist mit 9pt häufig zu klein
	% logofile=images/tuda_logo.pdf, %Falls die Logo Dateien nicht installiert sind
]{tudapub}

%%%%%%%%%%%%%%%%%%%%%%%%%%%%
% Download des TU-Logos
%%%%%%%%%%%%%%%%%%%%%%%%%%%%
% https://download.hrz.tu-darmstadt.de/protected/CE/TUDa_LaTeX/tuda_logo.pdf
% Der Pfad zum Logo kann als "logofile" angegeben werden.

%%%%%%%%%%%%%%%%%%%
% Sprachanpassung & Verbesserte Trennregeln
%%%%%%%%%%%%%%%%%%%
\usepackage[english, main=ngerman]{babel}
\usepackage[autostyle]{csquotes}% Anführungszeichen vereinfacht
\usepackage{microtype}

%%%%%%%%%%%%%%%%%%%
% Literaturverzeichnis
%%%%%%%%%%%%%%%%%%%
\usepackage[backend=biber,style=alphabetic]{biblatex}   % Literaturverzeichnis
\addbibresource{literature.bib}
% WICHTIG!! je nach Version ist es nötig, manuell auf Biber zum Erzeugen des Literaturverzeichnisses umzustellen: Optionen --> TexStudio konfigurieren... --> Erzeugen --> StandardBibliograhieprogramm --> Biber

%%%%%%%%%%%%%%%%%%%
% Paketvorschläge Tabellen
%%%%%%%%%%%%%%%%%%%
%\usepackage{array}     % Basispaket für Tabellenkonfiguration, wird von den folgenden automatisch geladen
\usepackage{tabularx}   % Tabellen, die sich automatisch der Breite anpassen
%\usepackage{longtable} % Mehrseitige Tabellen
%\usepackage{xltabular} % Mehrseitige Tabellen mit anpassarer Breite
\usepackage{booktabs}   % Verbesserte Möglichkeiten für Tabellenlayout über horizontale Linien

%%%%%%%%%%%%%%%%%%%
% Paketvorschläge Mathematik
%%%%%%%%%%%%%%%%%%%
\usepackage{mathtools} % erweiterte Fassung von amsmath
\usepackage{amssymb}   % erweiterter Zeichensatz
\usepackage{siunitx}   % Einheiten

%%%%%%%%%%%%%%%%%%%%%%
% Pseudocode
%%%%%%%%%%%%%%%%%%%
\usepackage[linesnumbered,lined,boxruled]{algorithm2e} % Package für Pseudocode

%%%%%%%%%%%%%%%%%%%
% Plotting und Grafik
%%%%%%%%%%%%%%%%%%%
\usepackage{tuda-pgfplots} % Package für Plotting with TUDa mods

%%%%%%%%%%%%%%%%%%%
% Sonstiges
%%%%%%%%%%%%%%%%%%%
\usepackage{blindtext} % Package für Blindtext

\begin{document}
	\title{WACE SoSe24 Hausarbeit Team03}
	\subtitle{Kurs: Wissenschaftliches Arbeiten im CE}
	%\author[V. Name]{Vorname Name} %optionales Argument ist die Signatur,
	\author[Karakoc, Rosenhagen, Tanase]{Aykut Karakoc, Christian Rosenhagen und Bogdan Tanase} %optionales Argument ist die Signatur,
	%\reviewer{Gutachter 1 \and Gutachterin 2} %Gutachten

	%Diese Felder erden untereinander auf der Titelseite platziert.
	%\department{ce} % Das Kürzel wird automatisch ersetzt und als Studienfach gewählt, siehe Liste der Kürzel im Dokument.
	\institution{
		% Fügt das CE-Logo über dem Schriftzug ein
		{\color{accentcolor}\includegraphics[width=0.9\linewidth]{images/ce_logo_black.pdf}}
		% oder etit-Logo: 
		%\includegraphics[width=0.9\linewidth]{images/etit_logo}
	}
	
	%\submissiondate{\today}
	%\examdate{\today}

	%	\tuprints{urn=1234,printid=12345}
	%	\dedication{Für alle, die \TeX{} nutzen.}

	\maketitle
	\pagenumbering{gobble} % Seitenzahlen angezeigt, startet ab dem Inhaltsverzeichnis

	%%%%%%%%%%%%%%%%%%%
	%Eigenständigkeitserklärung (bei Abschlussarbeiten BSc/MSc erforderlich, nicht hier bei dieser WACE-Hausarbeit!)
	%%%%%%%%%%%%%%%%%%%
	%\affidavit

	%%%%%%%%%%%%%%%%%%%
	%Abstract / Kurzzusammenfassung
	%%%%%%%%%%%%%%%%%%%


	%%%%%%%%%%%%%%%%%%%
	%Inhaltsverzeichnis 
	%%%%%%%%%%%%%%%%%%%
	\cleardoublepage
	\tableofcontents % Erstellte ein Inhaltsverzeichnis
	
	\cleardoublepage
	\pagenumbering{arabic} % Seitenzahlen angezeigt, startet ab dem Inhaltsverzeichnis
	\setcounter{page}{1} % Setzt den Seitenzahlenzähler auf 1

%%%%%%%%%%%%%%%%%%%%%%%%%%%%%%%%%%%%%%%%%%%%%%%%%%%%%%%%%%%%%%%%%%%%%%%%%%%%%%%%%%%%%%%%%%%%%%%%%%

	% INHALT, am Besten ausgelagert in eigene Files/Kapitel und dann mit \include{Unterordner/Filename} eingefügt, sorgt für bessere Übersichtlichkeit und Fehlersuche. Einzelne Dateien sind aktuell im Ordner Sections abgelegt. 
	
	
	%\include{}
	
	%\include{}
	
	%\include{}
	
	%\include{}
	
	%\include{}
	
%%%%%%%%%%%%%%%%%%%%%%%%%%%%%%%%%%%%%%%%%%%%%%%%%%%%%%%%%%%%%%%%%%%%%%%%%%%%%%%%%%%%%%%%%%%%%%%%%%
	
	%%%%%%%%%%%%%%%%%%%
	%Abbildungs- und Tabellenverzeichnis
	%%%%%%%%%%%%%%%%%%%
	%\listoffigures % Abbildungsverzeichnis (captions in den Figuren werden als Referenz genommen)
	%\listoftables % Verzeichnis der Tabellen (captions in den Tabellen werden als Referenz genommen)
	
	%%%%%%%%%%%%%%%%%%%
	%Literaturverzeichnis an dieser Stelle
	%%%%%%%%%%%%%%%%%%%
	%\printbibliography % Erstellt die Bibliography

\end{document}